\documentclass{controle}
\usepackage{main}

% \printanswers{}

\title{Évaluation n°12 : Vecteurs et algèbre}
\author{Seconde 3}
\date{Vendredi 6 Février}

\begin{document}
\titleandname{}
\version

\begin{questions}
\question[3]
Soit $\vect{u}$ un vecteur, et $k$ un nombre réel. Décrire la direction, le sens et la norme du vecteur $k\vect{u}$ par rapport à ceux du vecteur $\vect{u}$.

\ifprintanswers
\begin{itemize}
\item Direction : \textbf{Même direction que $\vect{u}$} 
\item Sens : \textbf{Si $k$ est positif, même sens que $\vect{u}$. Si $k$ est négatif, sens opposé à $\vect{u}$}
\item Norme : \textbf{On multiplie $\left|k\right|$ par la norme de $\vect{u}$.} 
\end{itemize}
\else
\begin{itemize}
\item Direction :
\item Sens : 
\item Norme :
\end{itemize}
\fi

\question[7]
On a placé trois points $A, B$ et $C$ sur le repère suivant :

\begin{center}
\begin{tikzpicture}
\draw[help lines] (-4.25,-4.25) grid[step=0.5] (4.25,4.25);
\coordinate (A) at (0,0);
\coordinate (B) at (0.5,1);
\coordinate (C) at (-1,1);
\draw (A) node[below] {$A$} node {$\times$};
\draw (B) node[above right] {$B$} node {$\times$};
\draw (C) node[above left] {$C$} node {$\times$};
\draw (A) -- (B) -- (C) -- cycle;

\ifprintanswers
\draw (A) ++ ($3*(A) - 3*(B)$) node {$\bullet$} node[below right] {$D$};
\draw (A) ++ ($(B) - (A)$) ++ ($3*(C) - 3*(A)$) node {$\bullet$} node[below right] {$E$};
\draw (C) ++ ($2*(C) - 2*(B)$) node {$\bullet$} node[below right] {$F$};
\else
\fi
\end{tikzpicture}  
\end{center}
Placer les points $D$, $E$ et $F$ pour que vérifier les égalités vectorielles suivantes.
\begin{parts}
\part[2] $\vect{AD} = 3\vect{BA}$
\part[2] $\vect{AE} = \vect{AB} - 3\vect{CA}$
\part[3] $3\vect{FC} - 2\vect{FB} = \vect{0}$ (Justifier la réponse) 
\end{parts}
\begin{solution}
On remarque l'égalité suivante :
\begin{equation*}
\begin{aligned}
3\vect{FC} - 2\vect{FB} &= 3\vect{FC} - 2(\vect{FC} + \vect{CB})&\text{~(Relation de Chasles)}\\
&= 3\vect{FC} -2\vect{FC} -2\vect{CB} &\text{~(Distributivité)}\\
&= \vect{FC} -2\vect{CB} &\text{~(Réduction)}
\end{aligned}
\end{equation*}
Ainsi, on cherche
\begin{equation*}
\begin{aligned}
&~\vect{FC} - 2\vect{CB} = \vect{0}\\
\Leftrightarrow& ~\vect{FC} = 2\vect{CB}\\
\Leftrightarrow& ~\vect{CF} = 2\vect{BC}\text{~(En multipliant chaque membre de l'égalité par $-1$)}
\end{aligned}
\end{equation*}
\end{solution}
\end{questions}
\newpage
\setcounter{page}{1}
\titleandname{}
\version
\begin{questions}
\question[3]
Soit $\vect{u}$ un vecteur, et $k$ un nombre réel. Décrire la direction, le sens et la norme du vecteur $k\vect{u}$ par rapport à ceux du vecteur $\vect{u}$.

\ifprintanswers
\begin{itemize}
\item Direction : \textbf{Même direction que $\vect{u}$} 
\item Sens : \textbf{Si $k$ est positif, même sens que $\vect{u}$. Si $k$ est négatif, sens opposé à $\vect{u}$}
\item Norme : \textbf{On multiplie $\left|k\right|$ par la norme de $\vect{u}$.} 
\end{itemize}
\else
\begin{itemize}
\item Direction :
\item Sens : 
\item Norme :
\end{itemize}
\fi

\question[7]
On a placé trois points $A, B$ et $C$ sur le repère suivant :

\begin{center}
\begin{tikzpicture}
\draw[help lines] (-4.25,-4.25) grid[step=0.5] (4.25,4.25);
\coordinate (A) at (0,0);
\coordinate (B) at (-1,0.5);
\coordinate (C) at (-0.5,-0.5);
\draw (A) node[above right] {$A$} node {$\times$};
\draw (B) node[above left] {$B$} node {$\times$};
\draw (C) node[below] {$C$} node {$\times$};
\draw (A) -- (B) -- (C) -- cycle;

\ifprintanswers
\draw (A) ++ ($3*(A) - 3*(B)$) node {$\bullet$} node[below right] {$D$};
\draw (A) ++ ($(B) - (A)$) ++ ($3*(C) - 3*(A)$) node {$\bullet$} node[below right] {$E$};
\draw (C) ++ ($2*(C) - 2*(B)$) node {$\bullet$} node[below right] {$F$};
\else
\fi
\end{tikzpicture}  
\end{center}
Placer les points $D$, $E$ et $F$ pour que vérifier les égalités vectorielles suivantes.
\begin{parts}
\part[2] $\vect{AD} = 3\vect{BA}$
\part[2] $\vect{AE} = \vect{AB} - 3\vect{CA}$
\part[3] $3\vect{FC} - 2\vect{FB} = \vect{0}$ (Justifier la réponse) 
\end{parts}
\begin{solution}
On remarque l'égalité suivante :
\begin{equation*}
\begin{aligned}
3\vect{FC} - 2\vect{FB} &= 3\vect{FC} - 2(\vect{FC} + \vect{CB})&\text{~(Relation de Chasles)}\\
&= 3\vect{FC} -2\vect{FC} -2\vect{CB} &\text{~(Distributivité)}\\
&= \vect{FC} -2\vect{CB} &\text{~(Réduction)}
\end{aligned}
\end{equation*}
Ainsi, on cherche
\begin{equation*}
\begin{aligned}
&~\vect{FC} - 2\vect{CB} = \vect{0}\\
\Leftrightarrow& ~\vect{FC} = 2\vect{CB}\\
\Leftrightarrow& ~\vect{CF} = 2\vect{BC}\text{~(En multipliant chaque membre de l'égalité par $-1$)}
\end{aligned}
\end{equation*}
\end{solution}
\end{questions}
\end{document}